% generated by benchmark_report.py - do not edit
H101-1 & System oznaczy brak ilości fenoksymetylpenicyliny (uwaga do recepty). & tak & \footnotesize Uwaga etapu 1: „fenoksymetylpenicylina: nie podano ilości”. \\
H101-2 & Moc „jeden milion jednostek” zostanie dopasowana do produktu 1 000 000 j.m. & \textbf{nie} & \footnotesize Model zapisał moc słownie („1 milion jednostek”) $\rightarrow$ brak produktu w rejestrze; produkt (Ospen 1000) wybrał lekarz. \\
H101-3 & Poza ilością leku projekty nie będą wymagały poprawek. & \textbf{nie} & \footnotesize Poza ilością potrzebne były poprawa mocy i wybór produktu oraz usunięcie dopisanego skrótu „(GAS)” w karcie. \\
H102-1 & Krople do oczu zostaną dopasowane do produktu w postaci kropli do oczu (nie innej postaci lewofloksacyny). & tak & \footnotesize Levofloxacin Genoptim 5 mg/ml, krople do oczu, roztwór. \\
H102-2 & Schemat malejący (co 2 godziny przez 2 dni, potem 4 razy dziennie) zostanie zachowany w dawkowaniu. & tak & \footnotesize Dawkowanie: „co dwie godziny w ciągu dnia przez pierwsze dwa dni, potem cztery razy dziennie, łącznie siedem dni”. \\
H102-3 & Zaprzeczenia (bez bólu, bez światłowstrętu) pozostaną zaprzeczeniami w karcie. & tak & \footnotesize Wywiad: „Bez bólu oka, bez światłowstrętu, ostrość widzenia niezaburzona”. \\
H103-1 & Recepta pozostanie pusta – paracetamol „który pacjent ma w domu” nie trafi na receptę. & tak & \footnotesize Recepta pusta; paracetamol tylko w zaleceniach („który masz w domu”). \\
H103-2 & Zaprzeczenia (bez duszności, bez krwioplucia, bez trzeszczeń) zostaną zachowane. & tak & \footnotesize Wywiad: „Bez duszności, bez bólu w klatce piersiowej, bez krwioplucia”; badanie: „bez trzeszczeń”. \\
H103-3 & Rozpoznanie zachowa niepewność („najpewniej wirusowe”). & tak & \footnotesize Rozpoznanie: „Ostre zapalenie oskrzeli, najpewniej wirusowe”. \\
H104-1 & Na recepcie znajdzie się tylko amlodypina 10 mg (atorwastatyna „bez zmian” nie). & tak & \footnotesize Na recepcie tylko amlodypina 10 mg (Adipine); atorwastatyna „bez zmian” w zaleceniach. \\
H104-2 & Zalecenie odstawienia dotychczasowej dawki 5 mg znajdzie się w zaleceniach. & tak & \footnotesize Zalecenie: „Dotychczasową amlodypinę 5 mg odstaw”. \\
H104-3 & Pomiary domowe i gabinetowe ciśnienia nie zostaną pomylone. & tak & \footnotesize Pomiary domowe 156/94 w wywiadzie, gabinetowe 158/96 i 152/94 w badaniu. \\
H105-1 & Moc „pięćdziesiąt mikrogramów” zostanie dopasowana do lewotyroksyny 50 mcg. & tak & \footnotesize „50 mikrogramów” $\rightarrow$ Eltroxin 50 mcg, tabletki. \\
H105-2 & Wartości laboratoryjne (TSH 12,4; FT4 0,7; anty-TPO 450) zostaną przepisane bez błędów. & tak & \footnotesize TSH 12,4 mIU/l, FT4 0,7 ng/dl, anty-TPO 450 – bez błędów. \\
H105-3 & Termin kontroli (TSH za 6–8 tygodni, potem wizyta) zostanie zachowany. & tak & \footnotesize „Badanie TSH za 6 do 8 tygodni, potem wizyta kontrolna z wynikiem”. \\
H106-1 & Metformina („zostawiamy, bez zmian”) nie trafi na receptę; empagliflozyna tak. & \textbf{nie} & \footnotesize Na recepcie tylko metformina („zostawiamy, bez zmian”), a empagliflozyny w ogóle nie było; lekarz zamienił pozycję (poprawka istotna). \\
H106-2 & Zasada odstawienia empagliflozyny przy wymiotach/biegunce znajdzie się w zaleceniach. & tak & \footnotesize „W razie wymiotów, biegunki albo gdy nie może pan jeść i pić, proszę na ten czas odstawić empagliflozynę”. \\
H106-3 & Wartości HbA1c (8,1\% obecnie, 7,6\% w marcu) nie zostaną zamienione. & tak & \footnotesize Wywiad: „HbA1c z ostatniego tygodnia 8,1\%, w marcu 7,6\%”. \\
H107-1 & Ryzyko: model pominie część przedłużanych leków, stosując regułę o lekach kontynuowanych bez zmian. & \textbf{nie} & \footnotesize Przewidywane ryzyko się nie zmaterializowało: wszystkie 5 przedłużanych leków na recepcie (2 wymagały wyboru postaci). \\
H107-2 & Ketoprofen (odstawiony) i paracetamol (tylko zalecenie) nie trafią na receptę. & tak & \footnotesize Ketoprofen jako zalecenie odstawienia, paracetamol w zaleceniach; żaden na recepcie. \\
H107-3 & Ilości opakowań pięciu leków zostaną przypisane do właściwych pozycji. & tak & \footnotesize Ilości (3$\times$28, 3$\times$30, 3$\times$28, 4$\times$60, 2$\times$50) przypisane do właściwych leków. \\
H108-1 & Symbicort (nazwa handlowa + budezonid z formoterolem) zostanie dopasowany do produktu 160 + 4,5 mcg/dawkę. & \textbf{nie} & \footnotesize Fałszywy brand\_substance\_conflict: substancje zapisane „budezonid z formoterolem”; produkt 160/4,5 wybrał lekarz. \\
H108-2 & Prednizon zostanie dopasowany do prednizonu, nie do prednizolonu. & tak & \footnotesize Prednizon 20 mg $\rightarrow$ Encorton 20 mg (prednizon). \\
H108-3 & Schemat „2 inhalacje 2 razy dziennie przez 2 tygodnie, potem 1” zostanie zachowany. & tak & \footnotesize „2 inhalacje 2 razy dziennie przez 2 tygodnie, potem 1 inhalacja 2 razy dziennie”. \\
H109-1 & Strona (prawa: Goldflam dodatni, ból) nie zostanie zamieniona z lewą (Goldflam ujemny). & tak & \footnotesize „Goldflama dodatni po stronie prawej, ujemny po stronie lewej”; rozpoznanie prawostronne. \\
H109-2 & Ujemny test ciążowy zostanie odnotowany w karcie. & tak & \footnotesize Wywiad: „Test ciążowy z moczu ujemny”. \\
H109-3 & Zalecenie o odstępie od wapnia, magnezu i żelaza zostanie zachowane. & tak & \footnotesize „… ani nie przyjmować go razem z preparatami wapnia, magnezu i żelaza. Zachować odstęp dwóch godzin”. \\
H110-1 & Recepta pozostanie pusta. & tak & \footnotesize Recepta pusta. \\
H110-2 & Obwody łydek (lewa 39 cm, prawa 36 cm) nie zostaną zamienione. & tak & \footnotesize „Obwód lewej łydki 39 cm, prawej 36 cm”. \\
H110-3 & Wstrzymanie HTZ znajdzie się w zaleceniach. & tak & \footnotesize „Hormonalną terapię zastępczą wstrzymaj do czasu wyjaśnienia sprawy”. \\
H110-4 & Termin kontroli („po wypisie z SOR”) nie zostanie pomylony ze skierowaniem na SOR. & tak & \footnotesize Skierowanie na SOR „dziś” w zaleceniach, kontrola osobno: „po wypisie z SOR, z dokumentacją”. \\
H111-1 & System oznaczy brak ilości hydroksyzyny. & tak & \footnotesize Uwaga etapu 1: „hydroksyzyna: nie podano ilości”. \\
H111-2 & Zaprzeczenie myśli samobójczych zostanie zachowane w karcie. & tak & \footnotesize Wywiad: „bez myśli rezygnacyjnych i samobójczych”. \\
H111-3 & Numer Centrum Wsparcia (800 70 2222) zostanie poprawnie złożony z dyktowania słownego. & tak & \footnotesize Numer złożony poprawnie: 800 70 2222 (razem z 112). \\
H112-1 & Ketoprofen w żelu 2,5\% zostanie dopasowany do produktu 25 mg/g w postaci żelu. & tak & \footnotesize „2,5\%” $\rightarrow$ Ketoprofen LGO 25 mg/g, żel. \\
H112-2 & Strona (prawe kolano) będzie spójna w recepcie, zaleceniach i karcie. & tak & \footnotesize Prawe kolano w dawkowaniu żelu, skierowaniu na RTG, wywiadzie, badaniu i rozpoznaniu. \\
H112-3 & Ostrzeżenie przed NLPZ doustnymi zostanie zachowane; żaden NLPZ doustny nie trafi na receptę. & tak & \footnotesize Ostrzeżenie przed doustnym ibuprofenem i ketoprofenem w zaleceniach; na recepcie tylko żel i paracetamol. \\
H113-1 & Xonvea zostanie dopasowana do produktu; kwas foliowy („kontynuować”) nie trafi na receptę. & \textbf{nie} & \footnotesize Częściowo: kwas foliowy poza receptą, ale Xonvea nie została wybrana automatycznie – fałszywy brand\_substance\_conflict („doksylamina z pirydoksyną”). \\
H113-2 & Warunkowe dawkowanie (2 tabletki wieczorem, ewentualnie dodatkowo 1 rano) zostanie zachowane. & tak & \footnotesize „na początek 2 tabletki wieczorem przed snem; jeśli po dwóch dniach nadal będą wymioty w ciągu dnia, dodatkowo 1 tabletka rano”. \\
H113-3 & Skala PUQE i ciała ketonowe zostaną odnotowane w karcie. & tak & \footnotesize Badanie: „ciała ketonowe śladowo. Skala PUQE: 9 pkt”. \\
H114-1 & Gastrolit nie zostanie wybrany automatycznie (produkt bez mocy w rejestrze) – system oznaczy wybór produktu. & tak & \footnotesize Uwaga: „Gastrolit: bez podanej mocy nie da się ustalić produktu - wybierz produkt”. \\
H114-2 & Paracetamol zostanie dopasowany do czopków 125 mg, nie do tabletek ani zawiesiny. & tak & \footnotesize Paracetamol Hasco 125 mg, czopki doodbytnicze. \\
H114-3 & Dwie masy ciała (12 kg obecnie, 12,5 kg tydzień temu) nie zostaną pomylone. & tak & \footnotesize „Masa ciała 12 kg; tydzień temu w domu 12,5 kg”. \\
H115-1 & „Amoksyklaw” (zapis fonetyczny) zostanie rozpoznany jako Amoksiklav 875 mg + 125 mg. & tak & \footnotesize „Amoksyklaw” $\rightarrow$ Amoksiklav 875 mg + 125 mg, tabletki powlekane. \\
H115-2 & Autokorekty zostaną uwzględnione: 10 dni i 2 opakowania (nie 7 dni i 1 opakowanie). & tak & \footnotesize Dawkowanie „przez 10 dni”, ilość „2 opakowania po 14 tabletek”. \\
H115-3 & Wtrącenia („yyy”, „eee”, „nie wiem jakie”) nie trafią do dokumentów. & tak & \footnotesize Wtrącenia nie trafiły do dokumentów („nie wiem jakie” $\rightarrow$ „nazwy nieznane”). \\
H115-4 & Mometazon zostanie dopasowany do aerozolu do nosa 50 mcg/dawkę. & tak & \footnotesize Mometazon Doppelherz 50 mcg/dawkę, aerozol do nosa, zawiesina. \\
H116-1 & System nie ostrzeże o sprzeczności amoksycyliny z alergią w karcie (brak takiej reguły) – wykryje to dopiero lekarz. & tak & \footnotesize Brak jakiejkolwiek uwagi systemu o alergii; amoksycylinę na klarytromycynę zamienił lekarz (poprawka istotna, wykrył lekarz). \\
H116-2 & Karta przepisze z dyktowania „nic nie uczula”, nie sygnalizując sprzeczności z danymi pacjenta. & \textbf{nie} & \footnotesize Karta nie przepisała samego „nic nie uczula”: zestawiła wprost zaprzeczenie pacjentki z alergią na amoksycylinę z karty, czyli sygnalizuje sprzeczność. \\
H116-3 & System oznaczy brak ilości amoksycyliny. & tak & \footnotesize Uwaga etapu 1: „amoksycylina: nie podano ilości”. \\
H117-1 & Bisoprolol („bez zmian, ma zapas”), ibuprofen (odstawiony) i paracetamol nie trafią na receptę. & tak & \footnotesize Na recepcie tylko furosemid, spironolakton i ramipril. \\
H117-2 & Ramipril („przedłużam”) trafi na receptę mimo reguły o lekach kontynuowanych. & tak & \footnotesize Ramipril 5 mg na recepcie (3 opakowania po 28). \\
H117-3 & Czasowe zwiększenie furosemidu (2 tabletki przez tydzień, potem 1) zostanie zachowane. & tak & \footnotesize „2 tabl. rano przez 7 dni, potem 1 tabl. rano”. \\
H117-4 & Brak postaci ramiprilu (kapsułki i tabletki w rejestrze) zostanie oznaczony do wyboru przez lekarza. & tak & \footnotesize Uwaga: „ramipril: lek występuje w tej mocy w kilku postaciach - wybierz produkt”. \\
H118-1 & W dawkowaniu amoksycyliny znajdzie się ostateczna objętość 10,5 ml, nie 10,8 ml z obliczeń. & tak & \footnotesize „10,5 ml dwa razy dziennie przez 7 dni”. \\
H118-2 & System oznaczy brak ilości ibuprofenu. & tak & \footnotesize Uwaga etapu 1: „ibuprofen: nie podano ilości”. \\
H118-3 & Strona (lewe płuco) nie zostanie zamieniona. & tak & \footnotesize „nad dolnym polem lewego płuca … prawe płuco bez zmian”; rozpoznanie lewostronne. \\
H118-4 & Amoksycylina zostanie dopasowana do zawiesiny 250 mg/5 ml. & tak & \footnotesize Ospamox 250 mg/5 ml, proszek do sporządzania zawiesiny doustnej. \\
H119-1 & Ibuprofen (odwołany) nie trafi na receptę, a zwolnienie lekarskie (odwołane) nie trafi do zaleceń. & tak & \footnotesize Ibuprofen nie na recepcie i nie w zaleceniach; zwolnienie lekarskie nie pojawia się w żadnym dokumencie. \\
H119-2 & Strona urazu (prawa) będzie poprawna w karcie mimo autokorekty w dyktowaniu. & tak & \footnotesize Prawy staw skokowy w wywiadzie, badaniu i rozpoznaniu; lewy „bez zmian”. \\
H119-3 & Voltaren Max bez podanej mocy nie zostanie wybrany automatycznie – system oznaczy wybór produktu. & tak & \footnotesize Uwaga: „diklofenak: bez podanej mocy nie da się ustalić produktu - wybierz produkt”. \\
H120-1 & Pilne skierowanie do okulisty „dziś” znajdzie się w zaleceniach jako pilne (nie jako zwykłe skierowanie). & tak & \footnotesize Pierwsze zalecenie: „Dziś pilnie zgłosić się na ostry dyżur okulistyczny zgodnie ze skierowaniem”. \\
H120-2 & Kontrola „za tydzień po okuliście” nie zostanie pomylona z pilnym skierowaniem. & tak & \footnotesize Kontrola osobno: „Kontrola u lekarza za tydzień, po wizycie u okulisty”. \\
H120-3 & Zaprzeczenia dotyczące oka (bez bólu, bez światłowstrętu) zostaną zachowane. & tak & \footnotesize Wywiad: „Brak bólu gałki ocznej. Brak światłowstrętu”. \\
